% Options for packages loaded elsewhere
\PassOptionsToPackage{unicode}{hyperref}
\PassOptionsToPackage{hyphens}{url}
\documentclass[smallcondensed]{svjour3}
\usepackage{xcolor}
\usepackage{amsmath,amssymb}
\setcounter{secnumdepth}{5}
\usepackage{iftex}
\ifPDFTeX
  \usepackage[T1]{fontenc}
  \usepackage[utf8]{inputenc}
  \usepackage{textcomp} % provide euro and other symbols
\else % if luatex or xetex
  \usepackage{unicode-math} % this also loads fontspec
  \defaultfontfeatures{Scale=MatchLowercase}
  \defaultfontfeatures[\rmfamily]{Ligatures=TeX,Scale=1}
\fi
\usepackage{lmodern}
\ifPDFTeX\else
  % xetex/luatex font selection
\fi
% Use upquote if available, for straight quotes in verbatim environments
\IfFileExists{upquote.sty}{\usepackage{upquote}}{}
\IfFileExists{microtype.sty}{% use microtype if available
  \usepackage[]{microtype}
  \UseMicrotypeSet[protrusion]{basicmath} % disable protrusion for tt fonts
}{}
\makeatletter
\@ifundefined{KOMAClassName}{% if non-KOMA class
  \IfFileExists{parskip.sty}{%
    \usepackage{parskip}
  }{% else
    \setlength{\parindent}{0pt}
    \setlength{\parskip}{6pt plus 2pt minus 1pt}}
}{% if KOMA class
  \KOMAoptions{parskip=half}}
\makeatother
\usepackage{color}
\usepackage{fancyvrb}
\newcommand{\VerbBar}{|}
\newcommand{\VERB}{\Verb[commandchars=\\\{\}]}
\DefineVerbatimEnvironment{Highlighting}{Verbatim}{commandchars=\\\{\}}
% Add ',fontsize=\small' for more characters per line
\newenvironment{Shaded}{}{}
\newcommand{\AlertTok}[1]{\textcolor[rgb]{1.00,0.00,0.00}{\textbf{#1}}}
\newcommand{\AnnotationTok}[1]{\textcolor[rgb]{0.38,0.63,0.69}{\textbf{\textit{#1}}}}
\newcommand{\AttributeTok}[1]{\textcolor[rgb]{0.49,0.56,0.16}{#1}}
\newcommand{\BaseNTok}[1]{\textcolor[rgb]{0.25,0.63,0.44}{#1}}
\newcommand{\BuiltInTok}[1]{\textcolor[rgb]{0.00,0.50,0.00}{#1}}
\newcommand{\CharTok}[1]{\textcolor[rgb]{0.25,0.44,0.63}{#1}}
\newcommand{\CommentTok}[1]{\textcolor[rgb]{0.38,0.63,0.69}{\textit{#1}}}
\newcommand{\CommentVarTok}[1]{\textcolor[rgb]{0.38,0.63,0.69}{\textbf{\textit{#1}}}}
\newcommand{\ConstantTok}[1]{\textcolor[rgb]{0.53,0.00,0.00}{#1}}
\newcommand{\ControlFlowTok}[1]{\textcolor[rgb]{0.00,0.44,0.13}{\textbf{#1}}}
\newcommand{\DataTypeTok}[1]{\textcolor[rgb]{0.56,0.13,0.00}{#1}}
\newcommand{\DecValTok}[1]{\textcolor[rgb]{0.25,0.63,0.44}{#1}}
\newcommand{\DocumentationTok}[1]{\textcolor[rgb]{0.73,0.13,0.13}{\textit{#1}}}
\newcommand{\ErrorTok}[1]{\textcolor[rgb]{1.00,0.00,0.00}{\textbf{#1}}}
\newcommand{\ExtensionTok}[1]{#1}
\newcommand{\FloatTok}[1]{\textcolor[rgb]{0.25,0.63,0.44}{#1}}
\newcommand{\FunctionTok}[1]{\textcolor[rgb]{0.02,0.16,0.49}{#1}}
\newcommand{\ImportTok}[1]{\textcolor[rgb]{0.00,0.50,0.00}{\textbf{#1}}}
\newcommand{\InformationTok}[1]{\textcolor[rgb]{0.38,0.63,0.69}{\textbf{\textit{#1}}}}
\newcommand{\KeywordTok}[1]{\textcolor[rgb]{0.00,0.44,0.13}{\textbf{#1}}}
\newcommand{\NormalTok}[1]{#1}
\newcommand{\OperatorTok}[1]{\textcolor[rgb]{0.40,0.40,0.40}{#1}}
\newcommand{\OtherTok}[1]{\textcolor[rgb]{0.00,0.44,0.13}{#1}}
\newcommand{\PreprocessorTok}[1]{\textcolor[rgb]{0.74,0.48,0.00}{#1}}
\newcommand{\RegionMarkerTok}[1]{#1}
\newcommand{\SpecialCharTok}[1]{\textcolor[rgb]{0.25,0.44,0.63}{#1}}
\newcommand{\SpecialStringTok}[1]{\textcolor[rgb]{0.73,0.40,0.53}{#1}}
\newcommand{\StringTok}[1]{\textcolor[rgb]{0.25,0.44,0.63}{#1}}
\newcommand{\VariableTok}[1]{\textcolor[rgb]{0.10,0.09,0.49}{#1}}
\newcommand{\VerbatimStringTok}[1]{\textcolor[rgb]{0.25,0.44,0.63}{#1}}
\newcommand{\WarningTok}[1]{\textcolor[rgb]{0.38,0.63,0.69}{\textbf{\textit{#1}}}}
\usepackage{longtable,booktabs,array}
\usepackage{calc} % for calculating minipage widths
% Correct order of tables after \paragraph or \subparagraph
\usepackage{etoolbox}
\makeatletter
\patchcmd\longtable{\par}{\if@noskipsec\mbox{}\fi\par}{}{}
\makeatother
% Allow footnotes in longtable head/foot
\IfFileExists{footnotehyper.sty}{\usepackage{footnotehyper}}{\usepackage{footnote}}
\makesavenoteenv{longtable}
\usepackage{graphicx}
\makeatletter
\newsavebox\pandoc@box
\newcommand*\pandocbounded[1]{% scales image to fit in text height/width
  \sbox\pandoc@box{#1}%
  \Gscale@div\@tempa{\textheight}{\dimexpr\ht\pandoc@box+\dp\pandoc@box\relax}%
  \Gscale@div\@tempb{\linewidth}{\wd\pandoc@box}%
  \ifdim\@tempb\p@<\@tempa\p@\let\@tempa\@tempb\fi% select the smaller of both
  \ifdim\@tempa\p@<\p@\scalebox{\@tempa}{\usebox\pandoc@box}%
  \else\usebox{\pandoc@box}%
  \fi%
}
% Set default figure placement to htbp
\def\fps@figure{htbp}
\makeatother
\setlength{\emergencystretch}{3em} % prevent overfull lines
\providecommand{\tightlist}{%
  \setlength{\itemsep}{0pt}\setlength{\parskip}{0pt}}
\journalname{Empirical Software Engineering}
\titlerunning{Posterior Guarantees Under Empirical Cost}
\authorrunning{Authors omitted}
\institute{Authors omitted for review}
\usepackage{float}
\usepackage{lineno}
\linenumbers
\setlength\linenumbersep{8pt}
\usepackage{bookmark}
\IfFileExists{xurl.sty}{\usepackage{xurl}}{} % add URL line breaks if available
\urlstyle{same}
\hypersetup{
  pdftitle={Posterior Guarantees Under Empirical Acquisition Cost: A Replication and Extension of FSE 2026},
  pdfauthor={Authors omitted for draft},
  hidelinks,
  pdfcreator={LaTeX via pandoc}}

\title{Posterior Guarantees Under Empirical Acquisition Cost: A
Replication and Extension of FSE 2026}
\author{Authors omitted for draft}
\date{}

\begin{document}
\maketitle

\begin{abstract}

Sampling is central to performance analysis of configurable software
systems: the full configuration space is often too large to measure
exhaustively, so researchers and practitioners select a finite sample
and reason from the observed behavior. A recent FSE 2026 paper proposed
a posterior-guarantee framework for this setting. Instead of asking only
whether a predictive model is accurate, it asks whether the worst
observed performance bound in a sample is trustworthy, using posterior
risk, posterior confidence, integrated confidence discrepancy, and
Statistical Recursive Search (SRS) {[}WMS26{]}. In parallel, recent work
on sampling cost argues that the usual proxy for acquisition effort,
sample size, is often not the same as the empirical cost of measuring
the sampled configurations {[}CostSamplingICSE{]}.

This paper connects these two lines of work. We first replicate the FSE
posterior-guarantee framework on the restored FSE artifact systems. We
then extend the analysis to x264, energy/performance measurements,
workload-sensitive configurable systems, Linux/TuxKConfig measured
pools, and LCBench measured hyperparameter-configuration pools. Across
936 primary context-budget scopes, posterior risk distinguishes sampling
strategies in 902 scopes, showing that the FSE metric transfers well.
However, replacing nominal sample size by measured acquisition cost
changes the confidence/cost Pareto frontier in 910 scopes. SRS remains
strong in the direct FSE revisit and especially in x264, but it is not a
universal empirical-cost-normalized recommendation: it is best in 54/102
x264 scopes, but 10/224 ENERGY scopes, 30/372 WORKLOAD scopes, 0/21
Linux scopes, and 0/210 LCBench scopes.

The contribution is therefore a refinement rather than a contradiction
of FSE 2026. Posterior risk and confidence are useful, portable
guarantee metrics, but the recommended sampling strategy depends on the
measured acquisition cost and the dataset context. Sample size remains a
useful nominal budget variable; it is not, by itself, a calibrated cost
model for posterior performance guarantees.

Keywords: configurable software systems, performance sampling, posterior
risk, posterior confidence, sampling cost, empirical cost, replication,
SRS.

\keywords{configurable software systems \and performance sampling \and posterior risk \and posterior confidence \and sampling cost \and empirical cost \and replication \and SRS}
\end{abstract}

\section{Introduction}\label{introduction}

Configurable software systems expose large configuration spaces. A
compiler, a database, a video encoder, or an operating-system kernel can
have tens, hundreds, or thousands of valid options, and the performance
impact of those options can vary by workload, input, and measured
property. Exhaustively measuring every configuration is often
impractical. The usual response is to sample: select a subset of
configurations, measure them, and use the sample to guide modeling,
optimization, or risk assessment.

Sampling creates a second problem. Once we have measured only a subset,
what can we say about the configurations that remain unmeasured? The FSE
2026 paper answers this question with a posterior-guarantee framework
{[}WMS26{]}. Given a sampled performance bound, such as the worst
performance value observed in the sample, the framework estimates the
posterior risk that unmeasured configurations violate that bound. It
also derives posterior confidence curves and an integrated confidence
discrepancy that compare the observed posterior guarantee with a prior
confidence baseline. The same paper proposes Statistical Recursive
Search (SRS), a sampling strategy designed to discover influential
decisions and improve the observed bound.

The FSE contribution is important because it changes the question from
``did our model predict well?'' to ``how much risk remains behind the
bound we observed?''. That question is directly relevant when a
practitioner uses sampling to make a safety, performance, or
resource-planning decision. If the sample never observes a slow
configuration, a large energy consumption, or a high-cost build, the
practitioner needs to know whether this absence is informative or merely
lucky.

At the same time, the sampling literature commonly treats sample size as
the cost of a study. If two strategies both measure n configurations,
they are often considered to have the same acquisition cost. This
assumption is convenient and sometimes defensible for abstract analysis,
but it is not generally true for performance measurement. Individual
configurations can differ substantially in runtime, energy use,
compilation time, storage footprint, or benchmarking duration. Two
samples with the same number of configurations can require very
different empirical effort. Recent cost-validity work makes this point
for accuracy-based sampling studies: sample size is a nominal budget,
not always a calibrated empirical cost {[}CostSamplingICSE{]}.

This paper combines these two ideas as a two-paper replication. From FSE
2026, we replicate posterior risk, posterior confidence, integrated
confidence discrepancy, and SRS. From the ICSE cost-validity study, we
replicate the construct-validity question: whether sample size is a safe
proxy for empirical sampling cost. We ask whether the FSE
posterior-guarantee results replicate and transfer, and whether the FSE
metrics lead to the same recommendations when empirical acquisition cost
replaces sample size. The goal is not to claim that the FSE framework is
wrong. On the contrary, our results show that posterior risk and
posterior confidence are useful across a wider set of systems and
datasets than the original paper. The goal is to test the cost model
behind the recommendation. If posterior guarantees are plotted, ranked,
or selected against sample size, do the same conclusions hold when the
x-axis is the measured cost actually paid to obtain the sample?

We study six primary corpora. The first is the restored FSE artifact,
which provides the closest direct replication of the original setting.
The second is x264 from the ICPE configurable-systems corpus, used as a
positive configurable-system replication with multiple inputs and target
properties {[}PereiraICPE20{]}. The third is an energy/performance
corpus, which changes the measured property used as acquisition cost
{[}TwinsEnergy{]}. The fourth is a workload corpus, which tests input
and workload sensitivity {[}WorkloadICSE23{]}. The fifth is a
Linux/TuxKConfig measured-pool stress test {[}TuxKConfig{]}. The sixth
is LCBench, a hyperparameter-configuration benchmark used as a
cross-domain measured-pool stress test {[}LCBench{]}.

Across these corpora, the main pattern is stable. Posterior risk
distinguishes strategies in almost all scopes. This supports the FSE
framework. However, empirical cost changes the confidence/cost Pareto
frontier in almost all scopes. This supports the cost-validity concern.
SRS remains a strong and often defensible strategy in the direct FSE
artifact and x264, but its advantage does not generalize uniformly to
ENERGY, WORKLOAD, Linux, or LCBench once empirical cost is used.

This paper makes four contributions:

\begin{enumerate}
\def\labelenumi{\arabic{enumi}.}
\tightlist
\item
  It provides a direct, artifact-based replication of the FSE
  posterior-risk framework on the restored FSE sampling data, including
  a documented comparison between the paper-text prior exponent and the
  artifact-compatible \texttt{n\ +\ 1} exponent.
\item
  It extends FSE-style posterior risk and confidence to five additional
  empirical settings: x264, ENERGY, WORKLOAD, Linux measured pools, and
  LCBench measured pools.
\item
  It evaluates the FSE metrics under measured empirical acquisition
  cost, showing that the cost-aware confidence Pareto frontier changes
  in 910/936 primary scopes.
\item
  It audits the local SRS implementation against the original FSE
  notebook cells on controlled fixtures and documents the one deliberate
  difference: the local implementation fills from unselected
  configurations so cost accounting does not double-count duplicate
  rows.
\end{enumerate}

The rest of the paper is organized as follows. Section 2 gives
background on posterior guarantees and empirical acquisition cost.
Section 3 states the research questions. Section 4 describes the
protocol, datasets, metrics, implementation, and statistical analysis.
Sections 5 and 6 present the results for replication/transfer and
empirical cost. Section 7 reports sensitivities and audits. Section 8
discusses implications. Section 9 covers threats to validity. Section 10
summarizes related work and Section 11 concludes.

\section{Background and Motivation}\label{background-and-motivation}

\subsection{Sampling for Configurable
Systems}\label{sampling-for-configurable-systems}

Configurable systems can expose a large number of valid configurations.
Each configuration is a vector of decisions, and each decision can
affect performance, energy, output size, compilation time, memory
footprint, or other non-functional properties. The full measurement
matrix is often unavailable: measuring every configuration under every
workload or input can be too slow, too expensive, or physically
impractical.

Sampling strategies choose which configurations to measure. Common
strategies include random sampling, pairwise or t-wise combinatorial
sampling, distance-based sampling, solver-based sampling,
diversity-based sampling, and adaptive strategies that use information
from previous measurements. The standard evaluation question is often
accuracy-based: for a fixed number of measurements, which strategy
supports the most accurate model or decision?

The posterior-guarantee perspective asks a different question. Suppose
the sample observes a worst-case bound, such as the maximum runtime in
the sampled configurations. A model might predict many values, but a
practitioner may simply need to know how credible the observed bound is.
If 100 configurations were measured and none exceeded a given bound, how
much risk remains that an unmeasured configuration does exceed it?

\subsection{FSE Posterior Risk and
Confidence}\label{fse-posterior-risk-and-confidence}

Let \texttt{C} be a finite reference set of configurations and let
\texttt{f(c)} be the measured property of configuration \texttt{c}. For
a max-bound analysis, a sample \texttt{S} observes the bound:

\begin{Shaded}
\begin{Highlighting}[]
\NormalTok{b(S) = max \{ f(c) | c in S \}.}
\end{Highlighting}
\end{Shaded}

The posterior risk of the sample is the fraction of configurations in
the reference set whose measured value exceeds that sampled bound:

\begin{Shaded}
\begin{Highlighting}[]
\NormalTok{risk(S) = | \{ c in C | f(c) \textgreater{} b(S) \} | / |C|.}
\end{Highlighting}
\end{Shaded}

For a min-bound analysis, the inequality is reversed. The present
primary protocol follows the FSE max-bound framing, with a documented
lower-bound sensitivity for throughput-like WORKLOAD cases.

Posterior confidence is the empirical distribution function of posterior
risks over repeated samples. For an accepted risk threshold
\texttt{epsilon}, posterior confidence is the fraction of repeated
samples whose posterior risk is at most \texttt{epsilon}:

\begin{Shaded}
\begin{Highlighting}[]
\NormalTok{confidence\_post(epsilon) = Pr[ risk(S) \textless{}= epsilon ].}
\end{Highlighting}
\end{Shaded}

The FSE framework compares this posterior confidence curve with a prior
confidence curve from scenario-theory reasoning. In the paper-first
protocol we use the paper-text exponent:

\begin{Shaded}
\begin{Highlighting}[]
\NormalTok{confidence\_prior(epsilon, n) = 1 {-} (1 {-} epsilon)\^{}n,}
\end{Highlighting}
\end{Shaded}

where \texttt{n} is the nominal sample size. The released artifact also
contains an \texttt{n\ +\ 1} convention; we rerun that convention as an
appendix sensitivity.

The integrated posterior confidence is the area under the posterior
confidence curve. For the empirical distribution used here, this is
exactly:

\begin{Shaded}
\begin{Highlighting}[]
\NormalTok{area\_post = 1 {-} mean(risk).}
\end{Highlighting}
\end{Shaded}

The integrated confidence discrepancy is the posterior area minus the
prior area. In the rerun tables this appears as
\texttt{omega\_size\_baseline} when the prior baseline uses nominal
sample size, and as \texttt{omega\_cost\_baseline} when the prior
baseline uses an effective sample size induced by measured acquisition
cost.

\subsection{SRS}\label{srs}

Statistical Recursive Search (SRS) is the adaptive strategy proposed in
the FSE artifact. It repeatedly draws a small random batch, identifies
decisions shared by high-performing configurations, tests whether those
decisions are statistically associated with performance, and recursively
restricts the search space. The artifact uses a default batch step size
of 25 and Welch's t-test at \texttt{p\ \textless{}\ 0.05}. If recursion
stops before the requested sample size, the sample is supplemented with
random configurations.

Our implementation follows these defaults in
\texttt{cost\_sampling/srs.py}. The direct FSE artifact rerun uses the
artifact-provided \texttt{SRSSampling\_t2} traces rather than
regenerating those samples. For generated or pool-based corpora, the
local SRS implementation generates samples with the same defaults and
100 seeds where the wrapper supports repeated stochastic runs.

\subsection{Sample Size as a Cost
Model}\label{sample-size-as-a-cost-model}

A sampling study needs a cost model. The simplest model is the count
model: one measured configuration costs one unit, so a sample of size
\texttt{n} costs \texttt{n}. This model is attractive because it is easy
to compare across systems, independent of hardware, and aligned with
many theoretical treatments.

The count model is not the same as empirical acquisition cost. If one
configuration takes 0.1 seconds to measure and another takes 100
seconds, a sample of 20 configurations can be cheap or expensive
depending on which configurations are selected. The same issue appears
when cost is energy, encoded output size, cumulative training time,
compilation time, or another measured property.

This distinction matters for posterior guarantees. A guarantee obtained
from 64 cheap configurations and a guarantee obtained from 64 expensive
configurations may have the same nominal sample size, but they did not
require the same measurement effort. A cost-aware decision should
therefore consider the measured acquisition cost behind the guarantee.

\subsection{Connecting the Two Lines of
Work}\label{connecting-the-two-lines-of-work}

The FSE framework supplies a useful outcome: posterior risk and
posterior confidence. The cost-validity argument supplies a different
x-axis: measured acquisition cost rather than nominal sample size.
Combining them asks:

\begin{Shaded}
\begin{Highlighting}[]
\NormalTok{Which strategy gives the best posterior guarantee for the empirical cost paid?}
\end{Highlighting}
\end{Shaded}

This question is not answered by either line of work alone. A
posterior-risk study can rank strategies by fixed sample size without
checking whether fixed sample size is fixed cost. A cost-validity study
can show that sample size is miscalibrated without asking whether the
resulting guarantee is trustworthy. The combined question is the focus
of this paper.

\section{Research Questions}\label{research-questions}

We organize the journal version around two main research questions.

\subsection{RQ1: Replication and Transfer of FSE Posterior
Claims}\label{rq1-replication-and-transfer-of-fse-posterior-claims}

\textbf{RQ1. To what extent do the FSE posterior-risk results replicate
on the original artifact systems and transfer to other datasets?}

This RQ evaluates the posterior-guarantee framework itself. It asks
whether posterior risk distinguishes strategies, whether guarantees
improve with nominal budget, whether SRS remains strong, and whether
these conclusions are context-sensitive.

We decompose RQ1 into four sub-questions:

\begin{itemize}
\tightlist
\item
  \textbf{RQ1.1 Metric transfer:} Do posterior-risk distributions
  distinguish sampling strategies at the same nominal budget?
\item
  \textbf{RQ1.2 Budget response:} Do posterior guarantees improve as the
  nominal sample budget grows?
\item
  \textbf{RQ1.3 SRS recommendation:} Does SRS remain a strong or best
  strategy under posterior-risk and posterior-confidence outcomes?
\item
  \textbf{RQ1.4 Context sensitivity:} Where do the answers change across
  systems, inputs, workloads, properties, and measured pools?
\end{itemize}

\subsection{RQ2: Replication of FSE Metrics Under Empirical
Cost}\label{rq2-replication-of-fse-metrics-under-empirical-cost}

\textbf{RQ2. Do FSE posterior metrics and recommendations replicate when
measured acquisition cost replaces sample size?}

This RQ evaluates the cost model behind the recommendation. It adapts
the existing cost-validity RQs to the posterior-guarantee setting.
Prediction error is replaced by posterior risk and integrated
confidence; sample-size cost is compared with measured empirical
acquisition cost.

We decompose RQ2 into five sub-questions:

\begin{itemize}
\tightlist
\item
  \textbf{RQ2.1 Cost heterogeneity:} At fixed sample size, do strategies
  induce different measured acquisition costs?
\item
  \textbf{RQ2.2 Sample-size proxy validity:} Is sample size a calibrated
  proxy for posterior-guarantee acquisition cost?
\item
  \textbf{RQ2.3 Recommendation stability:} Does the top guarantee/cost
  strategy change when empirical cost replaces sample size?
\item
  \textbf{RQ2.4 Pareto-front stability:} Does the confidence/cost Pareto
  frontier change when empirical cost replaces sample size?
\item
  \textbf{RQ2.5 Context sensitivity:} Are cost-aware conclusions stable
  across datasets, properties, inputs, workloads, and measured pools?
\end{itemize}

Table \ref{tab:replication-targets} makes the two replication targets
explicit. The first target provides the guarantee outcome and SRS
treatment. The second target provides the cost-model threat and the
sample-size-vs-measured-cost analysis pattern.

\begin{table}[H]
\caption{Replication targets and how they are used in this study.}
\label{tab:replication-targets}
\centering
\small
\begin{tabular}{@{}p{0.24\linewidth}p{0.69\linewidth}@{}}
\toprule
Target & Role in this paper \\
\midrule
FSE 2026 & Replicate posterior risk, posterior confidence, integrated discrepancy, SRS, and FSE-style claims; then test transfer beyond the released artifact. \\
ICSE cost study & Replicate the sample-size versus measured-cost threat model; replace prediction error with posterior guarantees. \\
Joint extension & Analyze sampling decisions in the measured guarantee/cost plane using aggregate counts, block tests, and curated frontier examples. \\
\bottomrule
\end{tabular}
\end{table}

\section{Experimental Protocol}\label{experimental-protocol}

\subsection{Units of Analysis}\label{units-of-analysis}

The main descriptive unit is a scope: one corpus-specific context at one
budget level. A context can be a system, an input/target pair, a
system/workload/property tuple, a Linux release, or an LCBench dataset,
depending on the corpus. A budget level is a nominal sample size such as
\texttt{t2}, \texttt{n32}, or a fraction of a measured Linux pool.

For each scope, we compute posterior risks over repeated runs of each
strategy. We then summarize:

\begin{itemize}
\tightlist
\item
  posterior-risk distributions by strategy;
\item
  median acquisition cost and percent of full measured cost;
\item
  effective sample size induced by measured cost;
\item
  integrated posterior confidence;
\item
  integrated confidence discrepancy against the size baseline;
\item
  integrated confidence discrepancy against the measured-cost baseline;
\item
  count-cost and measured-cost Pareto frontiers.
\end{itemize}

\subsection{Corpora}\label{corpora}

Table \ref{tab:corpora} summarizes the primary corpora and their
intended roles.

\begin{table}[H]
\caption{Primary corpora, budgets, measured cost axis, and replication role.}
\label{tab:corpora}
\centering
\scriptsize
\begin{tabular}{@{}p{0.23\linewidth}r p{0.12\linewidth}p{0.16\linewidth}p{0.30\linewidth}@{}}
\toprule
Corpus & Ctx. & Budgets & Cost & Role \\
\midrule
Original FSE artifact & 7 & t2 & performance & Direct FSE artifact revisit. \\
x264 & 34 & t1--t3 & time & Configurable-system replication. \\
ENERGY & 56 & n8--n64 & energy/perf. & Property/cost extension. \\
WORKLOAD & 93 & n8--n64 & time & Workload/input sensitivity. \\
Linux full & 7 & 1--5\% & vmlinux & Measured-pool stress test. \\
LCBench & 35 & n8--n64 & time\_total & HPO measured-pool stress test. \\
\bottomrule
\end{tabular}
\end{table}

The direct FSE artifact is the closest replication of the original
study. The official Zenodo \texttt{data.zip} was restored and
checksum-verified with SHA-256:

\begin{Shaded}
\begin{Highlighting}[]
\NormalTok{12ee2295d3d9d596eb57a92008b27562a0bf1882b21fcad9c8cd82082bf654bd}
\end{Highlighting}
\end{Shaded}

The Linux and LCBench studies are intentionally phrased as measured-pool
analyses. Their posterior risk is reference-relative to the measured
pool; we do not claim exhaustive risk over all possible Linux or
hyperparameter configurations.

\subsection{Strategies}\label{strategies}

The direct FSE artifact run uses the five paper-aligned strategies
available in the restored sampling traces:

\begin{itemize}
\tightlist
\item
  \texttt{random}
\item
  \texttt{solverBased}
\item
  \texttt{divDistBased}
\item
  \texttt{twise}
\item
  \texttt{SRSSampling}
\end{itemize}

The extended corpora use the strategy sets supported by their runners.
These include random, balanced, diversity, solver/distance or generated
baselines, pilot-cost strategies where available, and SRS. Diagnostic
cheap/expensive oracles are excluded from the primary generated-corpus
reruns because they are local diagnostic baselines rather than FSE paper
strategies.

\subsection{Posterior Metrics}\label{posterior-metrics}

For each repeated sample, we compute the sampled bound and posterior
risk over the reference universe. For each strategy and level, we
aggregate posterior risks over runs and compute integrated posterior
confidence:

\begin{Shaded}
\begin{Highlighting}[]
\NormalTok{area\_post = 1 {-} mean(risk).}
\end{Highlighting}
\end{Shaded}

We compare this area with two prior baselines. The size baseline uses
the nominal sample size \texttt{n}. The cost baseline converts the
median percent of full measured acquisition cost into an effective
sample size:

\begin{Shaded}
\begin{Highlighting}[]
\NormalTok{effective\_n\_by\_cost = (sample\_cost / full\_cost) * |C|.}
\end{Highlighting}
\end{Shaded}

The cost-aware confidence discrepancy is then:

\begin{Shaded}
\begin{Highlighting}[]
\NormalTok{omega\_cost = area\_post {-} area\_prior(effective\_n\_by\_cost).}
\end{Highlighting}
\end{Shaded}

This makes the cost-aware comparison structurally parallel to the
original size-aware comparison. It does not simply divide confidence by
time; it asks whether the posterior confidence exceeds the prior
confidence expected at the empirical cost actually paid.

\subsection{Decision Metrics}\label{decision-metrics}

For each scope, we compute six operational hypotheses:

\begin{longtable}[]{@{}lll@{}}
\toprule\noalign{}
Hyp. & Short question & Type \\
\midrule\noalign{}
\endhead
\bottomrule\noalign{}
\endlastfoot
FSE1 & Risk differs by strategy. & Inferential \\
FSE2 & Empirical cost changes the top choice. & Decision \\
FSE3 & SRS remains best under empirical cost. & Decision \\
FSE4 & Empirical cost changes the frontier. & Decision \\
FSE5 & SRS risk is lower than random. & Inferential \\
FSE6 & SRS and random costs differ. & Inferential \\
\end{longtable}

FSE1, FSE5, and FSE6 are inferential scope-level tests. FSE2, FSE3, and
FSE4 are deterministic decision indicators computed from the resulting
summaries and frontiers.

\subsection{Statistical Policy}\label{statistical-policy}

The statistical layer post-processes the clean rerun outputs without
modifying the raw posterior records. It applies Holm-Bonferroni
correction to the scope-level inferential hypotheses FSE1, FSE5, and
FSE6 within each corpus and hypothesis. Decision hypotheses FSE2, FSE3,
and FSE4 are reported as indicators, not p-value tests.

For paper-level claims, repeated budget levels and stochastic runs
should not be treated as independent systems. We therefore collapse
scope rows into corpus-specific blocks:

\begin{itemize}
\tightlist
\item
  original FSE artifact: system;
\item
  x264: input and target property;
\item
  ENERGY: system, target property, and cost property;
\item
  WORKLOAD: system, workload, target property, and cost property;
\item
  Linux: release;
\item
  LCBench: OpenML dataset.
\end{itemize}

For each corpus and hypothesis, the block majority supports or refutes
the hypothesis. Mixed blocks are reported separately and omitted from
the exact sign test denominator. Block-level p-values are Holm-corrected
within corpus across FSE1-FSE6.

\subsection{Protocol Constants and
Adaptations}\label{protocol-constants-and-adaptations}

The paper-first reruns use these constants:

\begin{itemize}
\tightlist
\item
  prior confidence exponent: paper-text \texttt{n}, implemented as
  \texttt{-\/-prior-offset\ 0};
\item
  SRS step size: 25;
\item
  SRS test: Welch t-test with \texttt{p\ \textless{}\ 0.05};
\item
  SRS completion: random fill to requested sample size;
\item
  repetitions: 100 stochastic/SRS seeds where the wrapper supports them;
\item
  primary diagnostic oracles: excluded;
\item
  primary bound orientation: max;
\item
  primary clean reruns: table-first, with curated plots generated post
  hoc.
\end{itemize}

The artifact-compatible \texttt{n\ +\ 1} prior exponent is retained as
an appendix sensitivity. Linux diversity sampling is the most expensive
part of the full Linux run; the completed full run uses two
256-candidate diversity runs per eligible release/level while keeping
random/balanced and SRS at 100 seeds. The earlier one-run, 128-candidate
Linux diversity pass is retained as a sampler budget sensitivity.

\section{RQ1 Results: Replication and
Transfer}\label{rq1-results-replication-and-transfer}

\subsection{Primary Descriptive
Results}\label{primary-descriptive-results}

Tables \ref{tab:rq1-descriptive} and \ref{tab:rq2-descriptive} summarize
the primary clean reruns. A scope is one context at one budget level.
The descriptive totals are useful as compact magnitude summaries, but
they should not be interpreted as pooled inferential tests because
corpora differ in size and structure.

\begin{table}[H]
\caption{RQ1-oriented descriptive outcomes under the clean rerun protocol.}
\label{tab:rq1-descriptive}
\centering
\scriptsize
\begin{tabular}{@{}lrrrrr@{}}
\toprule
Corpus & Scopes & Risk diff. & SRS best & SRS Pareto & SRS lower risk \\
\midrule
FSE artifact & 7 & 7/7 & 3/7 & 6/7 & 6/7 \\
x264 & 102 & 102/102 & 54/102 & 71/102 & 68/102 \\
ENERGY & 224 & 224/224 & 10/224 & 81/224 & 54/224 \\
WORKLOAD & 372 & 353/372 & 30/372 & 135/372 & 94/372 \\
Linux full & 21 & 21/21 & 0/21 & 15/21 & 5/21 \\
LCBench & 210 & 195/210 & 0/210 & 26/210 & 1/210 \\
Total & 936 & 902/936 & 97/936 & 334/936 & 228/936 \\
\bottomrule
\end{tabular}
\end{table}

\begin{table}[H]
\caption{RQ2-oriented descriptive outcomes when measured cost replaces sample count.}
\label{tab:rq2-descriptive}
\centering
\scriptsize
\begin{tabular}{@{}lrrrr@{}}
\toprule
Corpus & Scopes & Top changes & Frontier changes & SRS cost diff. \\
\midrule
FSE artifact & 7 & 2/7 & 6/7 & 7/7 \\
x264 & 102 & 7/102 & 91/102 & 66/102 \\
ENERGY & 224 & 51/224 & 224/224 & 106/224 \\
WORKLOAD & 372 & 72/372 & 359/372 & 156/372 \\
Linux full & 21 & 10/21 & 21/21 & 20/21 \\
LCBench & 210 & 133/210 & 209/210 & 26/210 \\
Total & 936 & 275/936 & 910/936 & 381/936 \\
\bottomrule
\end{tabular}
\end{table}

\begin{figure}
\centering
\pandocbounded{\includegraphics[keepaspectratio,alt={Aggregate posterior-guarantee and empirical-cost outcomes by corpus. The figure mirrors the table but emphasizes the contrast between metric transfer, SRS dominance, SRS Pareto membership, and frontier change.}]{paper/emse-svjour3/figures/fig_result_rates.png}}
\caption{Aggregate posterior-guarantee and empirical-cost outcomes by
corpus. The figure mirrors the table but emphasizes the contrast between
metric transfer, SRS dominance, SRS Pareto membership, and frontier
change.}
\end{figure}

\subsection{RQ1.1: Metric Transfer}\label{rq1.1-metric-transfer}

Posterior risk transfers well. Strategy posterior-risk distributions
differ in 902/936 primary scopes. The direct FSE artifact, x264, ENERGY,
and Linux show differences in every scope. WORKLOAD and LCBench are
slightly less uniform but still high at 353/372 and 195/210 scopes.

The block-level statistical layer supports the same conclusion for the
larger corpora. x264, ENERGY, WORKLOAD, and LCBench all support FSE1 at
the block level after correction. The direct FSE artifact and Linux each
have only seven blocks, so even all-directional scope results often
remain block-sign inconclusive after Holm correction. This should be
read as a power and blocking limitation, not as contradictory evidence.

\textbf{Answer to RQ1.1.} Yes. FSE-style posterior risk is a portable
outcome for comparing sampling strategies. It distinguishes strategies
across the direct artifact, configurable-system replications,
generated-sampling extensions, and measured-pool stress tests.

\subsection{RQ1.2: Budget Response}\label{rq1.2-budget-response}

Budget response is clearest in the corpora with comparable repeated
budget levels. x264 is the cleanest example: for SRS, posterior risk is
non-increasing in 34/34 input/target contexts and integrated posterior
confidence is non-decreasing in 34/34 contexts. Linux also shows
monotonic SRS posterior-risk improvement in the completed full pass.

This result should be interpreted carefully. A larger nominal sample
usually improves the posterior guarantee, but it does not imply that
sample size is a calibrated empirical cost. Two samples at the same
nominal level can have different measured costs, and a strategy can
improve posterior confidence while spending substantially more or less
acquisition effort.

\textbf{Answer to RQ1.2.} Mostly yes where budget scaling is directly
comparable. More samples generally improve posterior guarantees, but the
improvement should not be conflated with equal empirical cost.

\subsection{RQ1.3: SRS Recommendation}\label{rq1.3-srs-recommendation}

SRS remains important, but not universally dominant.

In the direct FSE artifact, SRS is best under the measured-cost baseline
in 3/7 systems and appears on the empirical cost/confidence Pareto
frontier in 6/7. This supports the original FSE intuition that SRS is a
strong strategy in the artifact setting.

x264 is the strongest positive external replication. SRS is best in
54/102 scopes and appears on the empirical Pareto frontier in 71/102. It
has lower posterior risk than random in 68/102 scopes. At the block
level, x264 supports the SRS-risk advantage over random and is
inconclusive, rather than negative, on the stricter ``SRS remains best
under empirical cost'' question.

The other corpora weaken a universal SRS recommendation. SRS is best in
only 10/224 ENERGY scopes and 30/372 WORKLOAD scopes. It is never the
best cost-normalized strategy in the Linux full run or LCBench. SRS
remains Pareto-efficient in many Linux scopes, but it is rarely
Pareto-efficient in LCBench.

Table \ref{tab:block-verdicts} compresses the block-level verdicts; the
heatmap figure shows the full corpus-by-hypothesis grid.

\begin{table}[H]
\caption{Block-level statistical verdicts, summarized by corpus.}
\label{tab:block-verdicts}
\centering
\scriptsize
\begin{tabular}{@{}p{0.24\linewidth}p{0.69\linewidth}@{}}
\toprule
Corpus & Block-level summary \\
\midrule
FSE artifact & Directional but inconclusive after blocking. \\
x264 & Supports risk transfer, SRS-risk advantage, cost differences, and frontier changes. \\
ENERGY & Supports risk transfer/frontier changes; refutes universal SRS best. \\
WORKLOAD & Supports risk transfer/frontier changes; refutes universal SRS best. \\
Linux full & Directional but inconclusive after seven release blocks. \\
LCBench & Supports risk transfer/frontier changes; refutes universal SRS best. \\
\bottomrule
\end{tabular}
\end{table}

The ``refutes'' label is hypothesis-specific. For FSE3, it means the
block-level evidence refutes the claim that SRS remains the best
empirical-cost-normalized strategy. For FSE2, it means top-choice
changes are not consistently frequent in that corpus. For FSE5, it means
the evidence does not support a broad SRS posterior-risk advantage over
random.

\begin{figure}
\centering
\pandocbounded{\includegraphics[keepaspectratio,alt={Block-level verdict heatmap for FSE1-FSE6. This figure follows the FSE-style hypothesis family while using ICSE-style blocking so repeated budgets do not become pseudo-replications.}]{paper/emse-svjour3/figures/fig_block_verdict_heatmap.png}}
\caption{Block-level verdict heatmap for FSE1-FSE6. This figure follows
the FSE-style hypothesis family while using ICSE-style blocking so
repeated budgets do not become pseudo-replications.}
\end{figure}

The direct FSE artifact deserves separate visual inspection because it
is the closest replication target. The next figure shows how much the
best measured-cost strategy improves over SRS in each original artifact
system. Green bars indicate systems where SRS is still the
cost-normalized winner; orange bars indicate systems where SRS remains
on the measured-cost frontier but is not best. An asterisk marks systems
where empirical cost changes the top recommendation.

\begin{figure}
\centering
\pandocbounded{\includegraphics[keepaspectratio,alt={Direct FSE artifact systems: measured-cost gap between the best strategy and SRS. This is the system-level companion to the aggregate SRS-best counts.}]{paper/emse-svjour3/figures/fig_fse_system_srs_gap.png}}
\caption{Direct FSE artifact systems: measured-cost gap between the best
strategy and SRS. This is the system-level companion to the aggregate
SRS-best counts.}
\end{figure}

\textbf{Answer to RQ1.3.} SRS remains a strong baseline and often a
defensible choice, especially for the direct FSE artifact and x264. It
should not be presented as a universally best cost-aware recommendation.

\subsection{RQ1.4: Context Sensitivity}\label{rq1.4-context-sensitivity}

The posterior framework transfers more robustly than the SRS
recommendation. x264 is favorable to SRS. ENERGY and WORKLOAD preserve
the usefulness of the posterior-risk metric but substantially reduce the
frequency with which SRS is the best empirical-cost-normalized strategy.
Linux shows SRS often remains Pareto-efficient but not top-ranked.
LCBench is the strongest negative stress-test case for SRS.

This is the central context-sensitivity result. The FSE posterior-risk
metric generalizes; a universal strategy recommendation does not.

The next figure gives FSE-style posterior-risk panels for three
representative artifact systems. These panels preserve the original
paper's visual language: they compare strategies by posterior risk at
the fixed artifact budget before the paper turns to the measured-cost
extension.

\begin{figure}
\centering
\pandocbounded{\includegraphics[keepaspectratio,alt={FSE-style posterior-risk panels for representative original artifact systems. These panels anchor the extension in the replicated FSE outcome before empirical acquisition cost is introduced.}]{paper/emse-svjour3/figures/fig_fse_key_posteriors.png}}
\caption{FSE-style posterior-risk panels for representative original
artifact systems. These panels anchor the extension in the replicated
FSE outcome before empirical acquisition cost is introduced.}
\end{figure}

\textbf{Answer to RQ1.4.} Context matters. The guarantee framework
transfers across contexts, but the strategy recommendation depends on
the corpus, workload, property, and acquisition-cost model.

\section{RQ2 Results: Empirical Acquisition
Cost}\label{rq2-results-empirical-acquisition-cost}

\subsection{RQ2.1: Cost Heterogeneity}\label{rq2.1-cost-heterogeneity}

At fixed nominal sample size, strategies often induce different measured
acquisition costs. SRS and random costs differ in 381/936 descriptive
scopes. The rate is especially high in the direct FSE artifact (7/7),
x264 (66/102), and Linux (20/21).

The lower rates in ENERGY, WORKLOAD, and LCBench should not be read as
evidence that sample size is generally safe. These corpora still show
extensive frontier changes, meaning that even when a particular pairwise
cost comparison is not consistently decisive, the overall guarantee/cost
decision surface can change.

\textbf{Answer to RQ2.1.} Fixed sample size does not generally imply
fixed measured acquisition cost. Cost heterogeneity is present and
strategy-dependent.

\subsection{RQ2.2: Sample-Size Proxy
Validity}\label{rq2.2-sample-size-proxy-validity}

The posterior rerun evaluates sample-size proxy validity by asking
whether the scientific decisions agree under two baselines:

\begin{itemize}
\tightlist
\item
  the size baseline, where prior confidence is computed using nominal
  \texttt{n};
\item
  the cost baseline, where prior confidence is computed using the
  effective sample size implied by measured acquisition cost.
\end{itemize}

If sample size were a safe proxy, the top recommendations and Pareto
frontiers would usually agree. They do not. The top recommendation
changes in 275/936 descriptive scopes, and the confidence/cost Pareto
frontier changes in 910/936.

Top changes are context-dependent; Pareto changes are nearly universal.
This is important. A top-1 winner is a scalar summary that can remain
stable even when many defensible tradeoff options change. The Pareto
frontier exposes the structure of the decision surface. Under empirical
acquisition cost, that surface is usually not the same as under nominal
sample size.

\textbf{Answer to RQ2.2.} Sample size is a useful nominal budget
parameter, but it is not a calibrated proxy for empirical acquisition
cost when making posterior-guarantee decisions.

\subsection{RQ2.3: Recommendation
Stability}\label{rq2.3-recommendation-stability}

The single top recommendation changes in some, but not all, corpora:

\begin{itemize}
\tightlist
\item
  direct FSE artifact: 2/7 scopes;
\item
  x264: 7/102 scopes;
\item
  ENERGY: 51/224 scopes;
\item
  WORKLOAD: 72/372 scopes;
\item
  Linux: 10/21 scopes;
\item
  LCBench: 133/210 scopes.
\end{itemize}

LCBench is the strongest top-change result. At the block level, LCBench
supports frequent top recommendation changes. x264, ENERGY, and WORKLOAD
refute frequent top changes at the block level, even though many
individual scopes change. The direct FSE artifact and Linux are
small-block and inconclusive.

The conclusion is therefore not that empirical cost always changes the
top choice. Instead, empirical cost sometimes changes the top choice,
and the frequency of that change is itself context-sensitive.

\textbf{Answer to RQ2.3.} Empirical cost can change the top
recommendation, but top-1 instability is not the most general result. It
is strongest in LCBench and visible in Linux and the direct artifact; it
is less frequent in x264, ENERGY, and WORKLOAD.

\subsection{RQ2.4: Pareto-Front
Stability}\label{rq2.4-pareto-front-stability}

Pareto-front stability is the strongest empirical-cost result. The
measured cost/confidence frontier changes in 910/936 descriptive scopes:

\begin{itemize}
\tightlist
\item
  direct FSE artifact: 6/7;
\item
  x264: 91/102;
\item
  ENERGY: 224/224;
\item
  WORKLOAD: 359/372;
\item
  Linux: 21/21;
\item
  LCBench: 209/210.
\end{itemize}

The block-level synthesis supports frontier changes for x264, ENERGY,
WORKLOAD, and LCBench. Direct FSE and Linux are directionally strong but
inconclusive at the corrected block-sign level because each has only
seven blocks.

This result is more important for practitioners than a top-1 flip alone.
A Pareto frontier represents the choices that are not obviously
wasteful: no other available strategy/level pair offers at least as much
confidence for no more cost. If empirical cost changes that set, then a
practitioner using sample-size cost may discard a defensible
measured-cost choice or retain a choice that is dominated once real
acquisition cost is considered.

\textbf{Answer to RQ2.4.} Yes. Replacing sample size with measured
acquisition cost changes the confidence/cost Pareto frontier in almost
every primary scope.

\subsection{RQ2.5: Cost-Aware Context
Sensitivity}\label{rq2.5-cost-aware-context-sensitivity}

The cost-aware conclusions differ by corpus.

x264 remains favorable to SRS: SRS is often best and often
Pareto-efficient. Nevertheless, empirical-cost frontiers change in
91/102 x264 scopes. This means the top recommendation can remain stable
while the surrounding set of defensible choices changes.

ENERGY and WORKLOAD show that the posterior metrics transfer, but SRS is
not often the best measured-cost choice. These corpora are important
because they remain within configurable-system style settings while
changing measured properties and workload contexts.

Linux shows a different pattern. SRS is never the best cost-normalized
strategy in the completed full run, but it remains on the measured-cost
Pareto frontier in 15/21 scopes. That is a nuanced negative result: SRS
is not dominant, but it is often not wasteful.

LCBench is the strongest external stress test. SRS is never the best
cost-normalized strategy and is Pareto-efficient in only 26/210 scopes.
This does not refute FSE for configurable software; it shows that the
SRS recommendation should not be generalized without checking domain and
cost structure.

\textbf{Answer to RQ2.5.} The effect of empirical cost is
context-dependent. The most stable result is not a universal winner
change, but a widespread change in the guarantee/cost Pareto frontier.

\subsection{Curated Frontier Examples}\label{curated-frontier-examples}

The curated frontier example set consists of five panels generated by
the curated-frontier artifact. These panels illustrate the main decision
patterns without dumping every per-context plot.

\begin{figure}
\centering
\pandocbounded{\includegraphics[keepaspectratio,alt={Summary of the five curated frontier examples. The bars show the best-minus-SRS omega-cost gap; green means SRS is best, orange means SRS is on the frontier but not best, and red means SRS is not on the measured-cost frontier.}]{paper/emse-svjour3/figures/fig_curated_srs_gap.png}}
\caption{Summary of the five curated frontier examples. The bars show
the best-minus-SRS omega-cost gap; green means SRS is best, orange means
SRS is on the frontier but not best, and red means SRS is not on the
measured-cost frontier.}
\end{figure}

\begin{figure}
\centering
\pandocbounded{\includegraphics[keepaspectratio,alt={Direct FSE 7z frontier}]{artifacts/fse_posterior_curated_frontiers/01_fse_7z_srs_pareto_not_best.png}}
\caption{Direct FSE 7z frontier}
\end{figure}

Panel A: the restored FSE 7z system. SRS remains Pareto-efficient, but
measured cost changes the top recommendation to \texttt{twise\_t2}.

\begin{figure}
\centering
\pandocbounded{\includegraphics[keepaspectratio,alt={x264 frontier}]{artifacts/fse_posterior_curated_frontiers/02_x264_time_1_srs_best.png}}
\caption{x264 frontier}
\end{figure}

Panel B: x264 time input 1. SRS is both the measured-cost winner and
Pareto-efficient.

\begin{figure}
\centering
\pandocbounded{\includegraphics[keepaspectratio,alt={WORKLOAD frontier}]{artifacts/fse_posterior_curated_frontiers/03_workload_kanzi_vmlinux_srs_pareto_not_best.png}}
\caption{WORKLOAD frontier}
\end{figure}

Panel C: WORKLOAD kanzi/vmlinux. SRS is Pareto-efficient but not the
best measured-cost option.

\begin{figure}
\centering
\pandocbounded{\includegraphics[keepaspectratio,alt={Linux frontier}]{artifacts/fse_posterior_curated_frontiers/04_linux500_pool_srs_pareto_not_best.png}}
\caption{Linux frontier}
\end{figure}

Panel D: Linux 5.0.0 measured pool. SRS is Pareto-efficient but not
dominant.

\begin{figure}
\centering
\pandocbounded{\includegraphics[keepaspectratio,alt={LCBench frontier}]{artifacts/fse_posterior_curated_frontiers/05_lcbench_mfeat_factors_srs_not_frontier.png}}
\caption{LCBench frontier}
\end{figure}

Panel E: LCBench mfeat-factors. SRS is not on the measured-cost
frontier.

For journal layout, these panels can be combined into one multi-panel
figure with a shared legend and a caption explaining that the x-axis is
measured acquisition cost and the y-axis is integrated posterior
confidence.

\section{Sensitivities and Audits}\label{sensitivities-and-audits}

\subsection{Prior Exponent
Sensitivity}\label{prior-exponent-sensitivity}

The paper-first protocol uses the paper-text prior exponent \texttt{n}.
The released artifact convention uses \texttt{n\ +\ 1}. We reran the
direct FSE artifact with the artifact-compatible exponent as an appendix
sensitivity.

The headline counts are identical:

\begin{itemize}
\tightlist
\item
  SRS best: 3/7;
\item
  SRS Pareto: 6/7;
\item
  top recommendation changes: 2/7;
\item
  frontier changes: 6/7.
\end{itemize}

This means the main conclusion is not driven by the exponent convention.
The \texttt{n\ +\ 1} run should remain in the appendix, because the
paper-text definition is the clearer main-text choice and the
sensitivity validates robustness.

\subsection{x264 Size as Cost}\label{x264-size-as-cost}

The x264 primary run uses encoding time as acquisition cost while
evaluating time and size targets. We also reran x264 with encoded size
as acquisition cost. SRS remains favorable:

\begin{itemize}
\tightlist
\item
  SRS best: 57/102;
\item
  SRS Pareto: 76/102;
\item
  frontier changes: 96/102.
\end{itemize}

This supports the claim that x264 is a positive external replication for
SRS under both primary cost choices, while still preserving the
cost-aware frontier result.

\subsection{WORKLOAD Throughput
Framing}\label{workload-throughput-framing}

The primary WORKLOAD run treats \texttt{h2} throughput as an inverse
acquisition-cost quantity. We also ran an \texttt{h2} lower-bound
framing sensitivity. In that sensitivity:

\begin{itemize}
\tightlist
\item
  SRS best: 1/32;
\item
  SRS Pareto: 5/32;
\item
  frontier changes: 31/32.
\end{itemize}

This sensitivity does not make SRS dominant. It supports moving detailed
throughput-orientation discussion to the appendix.

\subsection{LCBench Cumulative
Runtime}\label{lcbench-cumulative-runtime}

The primary LCBench run uses the last observed runtime mode. A
cumulative runtime sensitivity gives the same headline counts:

\begin{itemize}
\tightlist
\item
  SRS best: 0/210;
\item
  SRS Pareto: 26/210;
\item
  frontier changes: 209/210.
\end{itemize}

The LCBench stress-test conclusion is robust to this runtime-mode
choice.

\subsection{Linux Diversity Budget}\label{linux-diversity-budget}

The completed Linux full run uses two 256-candidate diversity runs per
eligible release/level. An earlier one-run, 128-candidate diversity pass
is retained as a sampler-budget sensitivity. It preserves the key
conclusions:

\begin{itemize}
\tightlist
\item
  SRS best: 0/21;
\item
  SRS Pareto: 15/21;
\item
  frontier changes: 21/21.
\end{itemize}

This suggests that increasing the diversity candidate budget strengthens
some inferential details but does not change the SRS-best or Pareto
conclusions.

\subsection{SRS Notebook Replay Audit}\label{srs-notebook-replay-audit}

The local SRS implementation is a code-native implementation aligned
with the FSE artifact notebook. The preserved notebook snapshot is:

\begin{Shaded}
\begin{Highlighting}[]
\NormalTok{artifacts/srs\_original\_fse\_implementation/SRS{-}sampling.ipynb}
\NormalTok{SHA{-}256: 24769e71e5827a98278cbb148ba1268d4d446d11fb5816fe62955d3994e250f0}
\end{Highlighting}
\end{Shaded}

The notebook replay audit executes the preserved notebook import/helper
cells on controlled fixtures and compares row outputs with
\texttt{cost\_sampling.srs}. Three cases are recorded:

\begin{itemize}
\tightlist
\item
  a single-batch fixture with exact row identity match;
\item
  a recursive strong-signal fixture with exact row identity match;
\item
  a duplicate-fill fixture with the expected documented divergence.
\end{itemize}

The divergence is intentional. The notebook may fill from the original
measurement pool and repeat rows already selected. The local
implementation fills from remaining rows so acquisition-cost accounting
does not double-count a configuration. For strict bit-for-bit historical
replay, the audit harness should be run in the original artifact's
pandas/SciPy environment. For the present paper, the correct wording is
``faithful code-native implementation with controlled output-level
replay,'' not ``bit-for-bit artifact reuse under all library versions.''

\section{Discussion}\label{discussion}

\subsection{A Replication, Not a
Rejection}\label{a-replication-not-a-rejection}

The results should be read as a replication-plus-extension of FSE 2026.
The posterior-risk metric works. It distinguishes strategies across many
corpora, and it gives a meaningful guarantee-oriented view that is
different from prediction error. The direct artifact revisit also shows
that SRS remains strong in the original setting.

The extension is about cost. The FSE recommendation is naturally stated
at a fixed nominal sample size. We ask whether the same recommendation
holds when the acquisition cost is the measured effort required to
obtain the sample. The answer is nuanced: often yes for x264, sometimes
for the direct artifact, but not generally across all external settings.

\subsection{Why Pareto Frontiers
Matter}\label{why-pareto-frontiers-matter}

Top-1 recommendations are easy to report but brittle. A single scalar
winner depends on the objective, normalization, and ties. Pareto
frontiers are more decision-oriented. They show which choices are not
dominated in the confidence/cost plane.

The strongest empirical-cost result is that Pareto frontiers change in
910/936 primary scopes. This means empirical cost changes the decision
set even when the single top choice remains stable. For practitioners,
this is more useful than simply asking whether SRS wins. It identifies
which strategy/budget pairs remain defensible once the actual
measurement cost is considered.

\subsection{SRS as a Strong Baseline, Not a Universal
Winner}\label{srs-as-a-strong-baseline-not-a-universal-winner}

The paper should avoid overstating the negative result. SRS is not
``bad.'' It is often Pareto-efficient and sometimes best. In x264, it is
the strongest positive external result. In Linux, it is never the top
cost-normalized choice but remains Pareto-efficient in most scopes. In
the direct FSE artifact, it is Pareto-efficient in 6/7 systems.

The correct conclusion is that SRS should be treated as a strong
baseline whose cost-aware value depends on the dataset and cost axis. A
practitioner should not assume that SRS is the best empirical-cost
strategy without measuring the cost profile of the candidate samples.

\subsection{Effective Cost Baselines Are More Informative Than Raw
Ratios
Alone}\label{effective-cost-baselines-are-more-informative-than-raw-ratios-alone}

The cost-aware posterior metric used here is not merely ``confidence
divided by cost.'' It converts measured acquisition cost into an
effective sample size and then asks whether the observed posterior
confidence exceeds the prior confidence expected at that effective
budget. This keeps the comparison close to the FSE framework while
replacing the cost abstraction.

Raw calibration ratios remain useful for diagnosing cost heterogeneity,
but the posterior question is ultimately decision-oriented: which
choices remain defensible in the confidence/cost plane?

\subsection{Practical Implications}\label{practical-implications}

For researchers, the implication is methodological. Reporting posterior
guarantees only at fixed sample sizes can hide cost-model assumptions.
Future sampling studies should report both nominal size and measured
acquisition cost when the measured property is heterogeneous.

For practitioners, the implication is operational. A sampling strategy
should be selected not only by how many configurations it measures, but
by the guarantee obtained for the measured cost spent. A strategy that
samples fewer configurations is not necessarily cheaper; a strategy that
samples the same number of configurations is not necessarily equal cost.

For artifact evaluators, the implication is reproducibility-oriented.
The direct FSE rerun required restoring the official artifact data and
documenting the prior-exponent convention. The SRS audit required
preserving the original notebook and comparing outputs on controlled
fixtures. These details matter when a replication shifts from metric
reproduction to cost-aware extension.

\section{Threats to Validity}\label{threats-to-validity}

\subsection{Construct Validity}\label{construct-validity}

The main construct threat is the definition of acquisition cost.
Different corpora expose different measured properties. In x264, primary
acquisition cost is encoding time and appendix cost is encoded size. In
ENERGY, cost can be energy or performance. In WORKLOAD, cost is time
except for throughput-like cases. In Linux, the measured pool uses
\texttt{vmlinux} as the cost-relevant quantity. In LCBench, cost is
\texttt{time\_total}.

This heterogeneity is intentional because the research question is about
measured acquisition cost, not one universal cost unit. Still, the final
paper must be explicit that cost is corpus-specific and that
cross-corpus totals are descriptive, not a single pooled cost scale.

A second construct threat is the use of measured-pool posterior risk for
Linux and LCBench. These are not exhaustive configuration spaces. We
address this by using ``pool-relative'' or ``reference-relative''
posterior risk for those corpora.

\subsection{Internal Validity}\label{internal-validity}

The posterior implementation was written locally and aligned with the
FSE artifact formulas. The prior exponent differs between the paper text
and the artifact notebook; we make the paper-text \texttt{n} convention
primary and rerun the artifact-compatible \texttt{n\ +\ 1} convention as
a sensitivity. The sensitivity does not change the direct artifact
headline counts.

The SRS implementation is code-native. Controlled notebook replay shows
exact matches on two fixtures and a documented duplicate-fill divergence
on one fixture. The direct FSE artifact run uses the artifact's shipped
SRS traces, reducing risk for the direct replication. Generated and
measured-pool corpora use the local implementation.

Linux diversity sampling required a runtime adaptation. The completed
full run uses two 256-candidate diversity runs per eligible level; the
earlier one-run 128-candidate pass is retained as a sensitivity and
preserves the key SRS-best and Pareto conclusions.

\subsection{Statistical Conclusion
Validity}\label{statistical-conclusion-validity}

Scope-level rows are not all independent. A corpus can contribute many
budget levels, inputs, or repeated stochastic runs. We therefore use
block-level sign tests for paper-level claims, with blocks defined by
systems, input/target pairs, workload contexts, releases, or datasets as
appropriate.

The block-level tests are conservative for small-block corpora. The
direct FSE artifact and Linux each have only seven blocks, so strong
descriptive counts can remain statistically inconclusive after Holm
correction. We report both scope-level magnitudes and block-level
verdicts to avoid overclaiming.

\subsection{External Validity}\label{external-validity}

The six corpora are broad but not exhaustive. The direct FSE artifact
covers the original setting. x264, ENERGY, and WORKLOAD broaden
configurable-system evidence. Linux and LCBench provide measured-pool
stress tests. Other systems, hardware platforms, workloads, and
measurement protocols may yield different cost structures.

The result should therefore be generalized as a methodological caution:
posterior guarantees should be evaluated against measured cost when
acquisition cost is heterogeneous. It should not be generalized as a
claim that one particular alternative strategy always dominates SRS.

\subsection{Reproducibility}\label{reproducibility}

All primary and appendix reruns are recorded as artifact directories in
this repository. The clean protocol, compact synthesis, statistical
tests, curated frontiers, SRS notebook audit, and implementation modules
are part of the artifact package.

The restored FSE data archive is checksum-verified, but some raw
external inputs, such as LCBench JSON and workload archives, remain
local external data dependencies. Final artifact packaging should either
provide those inputs, document download instructions, or provide
checksums and expected paths.

\section{Related Work}\label{related-work}

\subsection{Posterior Guarantees for Performance
Sampling}\label{posterior-guarantees-for-performance-sampling}

The closest related work is the FSE 2026 posterior-risk paper that
introduced the guarantee framework and SRS {[}WMS26{]}. Our work
directly replicates that framework on the restored artifact and then
extends it to measured empirical acquisition cost. The present study
should therefore be positioned as a replication and cost-aware
extension, not as an unrelated competing method.

\subsection{Sampling Strategies for Configurable
Systems}\label{sampling-strategies-for-configurable-systems}

Configurable-system sampling has studied random, combinatorial,
distance-based, solver-based, diversity-based, and adaptive strategies
{[}KalteneckerICSE19, PereiraICPE20{]}. Much of this work uses sample
size as the primary cost axis. The x264 and broader configurable-system
corpora used here inherit that tradition, but we reinterpret their
samples through measured cost and posterior guarantees.

\subsection{Empirical Cost of
Sampling}\label{empirical-cost-of-sampling}

The cost-validity line of work argues that sample size is not
necessarily sampling cost {[}CostSamplingICSE{]}. It studies whether
conclusions based on count-cost differ from conclusions based on
measured acquisition time, energy, or related properties. Earlier
cost-aware configurable-systems work also considered measurement effort,
but often through count-based abstractions such as fixed train/test
acquisition counts {[}SarkarASE15{]}. This paper reuses the
empirical-cost framing but changes the outcome from prediction error to
posterior risk and confidence.

\subsection{Workload, Input, and Property
Sensitivity}\label{workload-input-and-property-sensitivity}

Prior configurable-system studies show that performance behavior can
depend on inputs, workloads, and non-functional properties
{[}PereiraICPE20, WorkloadICSE23, TwinsEnergy{]}. Our WORKLOAD and
ENERGY extensions are aligned with that literature. The novelty here is
the cost-aware posterior-guarantee question: not merely whether behavior
changes across contexts, but whether the posterior guarantee/cost
decision surface changes.

\subsection{Hyperparameter and Measured-Pool
Benchmarks}\label{hyperparameter-and-measured-pool-benchmarks}

LCBench and Linux/TuxKConfig measured pools are not direct replicas of
the FSE configuration-space setting {[}LCBench, TuxKConfig{]}. They are
stress tests for the methodology. They ask whether the
posterior-guarantee and empirical-cost reasoning behaves similarly in
large measured pools with heterogeneous acquisition costs. This helps
delimit the scope of the SRS recommendation.

\section{Artifact and Data
Availability}\label{artifact-and-data-availability}

The current draft is backed by the following committed artifact groups:

\begin{itemize}
\tightlist
\item
  protocol and synthesis reports;
\item
  statistical-testing outputs;
\item
  curated frontier figures and manifest;
\item
  SRS notebook audit and notebook replay notes;
\item
  SRS and posterior implementation modules.
\end{itemize}

Primary rerun directories:

\begin{itemize}
\tightlist
\item
  direct FSE artifact clean paper protocol;
\item
  x264 ICPE paper protocol;
\item
  ENERGY paper protocol;
\item
  WORKLOAD paper protocol;
\item
  Linux full diversity protocol;
\item
  LCBench 100-seed paper protocol.
\end{itemize}

Appendix rerun directories:

\begin{itemize}
\tightlist
\item
  direct FSE artifact \texttt{n\ +\ 1} sensitivity;
\item
  x264 size-cost sensitivity;
\item
  WORKLOAD throughput sensitivity;
\item
  LCBench cumulative-runtime sensitivity;
\item
  Linux full default-diversity sensitivity.
\end{itemize}

Before submission, this section should be converted into the journal's
artifact availability format, with persistent archive identifiers and
checksums for all external raw inputs.

\section{Conclusion}\label{conclusion}

This paper revisits FSE 2026 posterior performance guarantees under a
measured acquisition-cost model. The result is a refinement of the FSE
story. Posterior risk and posterior confidence replicate and transfer
well: they distinguish sampling strategies in 902/936 primary scopes.
SRS remains a strong strategy in the direct FSE artifact and x264.
However, sample size is not a reliable proxy for empirical acquisition
cost when making guarantee/cost decisions. Replacing nominal sample size
with measured cost changes the confidence/cost Pareto frontier in
910/936 primary scopes, and SRS is not the universal cost-normalized
winner across ENERGY, WORKLOAD, Linux, and LCBench.

The practical message is straightforward. Posterior guarantees answer an
important question: how much risk remains behind a sampled bound? To
make that answer actionable, the guarantee should be interpreted against
the empirical cost paid to obtain it. Sample size remains useful for
experiment design, but measured acquisition cost is the relevant budget
for cost-aware decisions.

\section{Draft Reference Keys}\label{draft-reference-keys}

The current repository does not contain a BibTeX file. The following
draft keys identify the references that need bibliographic metadata
before LaTeX conversion or journal submission.

\begin{itemize}
\tightlist
\item
  \texttt{{[}WMS26{]}} FSE 2026 posterior-risk and SRS paper under
  replication.
\item
  \texttt{{[}CostSamplingICSE{]}} ``Sample Size Is Not Sampling Cost''
  ICSE submission and public reproducibility repository.
\item
  \texttt{{[}KalteneckerICSE19{]}} Configurable-system sampling corpus
  used in the REPLICATION family.
\item
  \texttt{{[}PereiraICPE20{]}} x264 ICPE configurable-system artifact
  and dataset.
\item
  \texttt{{[}TwinsEnergy{]}} Twins/Twins-or-False-Friends
  energy/performance corpus.
\item
  \texttt{{[}WorkloadICSE23{]}} Workload-impact configurable-system
  corpus.
\item
  \texttt{{[}TuxKConfig{]}} Linux/TuxKConfig measured-pool data.
\item
  \texttt{{[}LCBench{]}} LCBench lightweight benchmark data.
\item
  \texttt{{[}SarkarASE15{]}} Prior cost-aware configurable-systems
  sampling work using a count-based cost abstraction.
\end{itemize}

\end{document}
